%% Section 6: Analysis and Discussion
%% English manuscript section.
%% -----------------------------------------------------------------------

\section{Analysis and Discussion}
\label{sec:analysis}

\subsection{Ablation Study: Contribution of the Modulo Stream}
\label{sec:ablation}

To isolate the contribution of the modulo stream and FiLM fusion, we compare IntSeqBERT with
the magnitude-only Ablation model. As seen in Table~\ref{tab:main}, the modulo stream
delivers the largest gains in modulo prediction (MMA $+15.2$\,pt, parity $+13.5$\,pt at Large
scale). Its contribution to magnitude prediction is smaller but significant (Acc$_{0.5}$
$+6.2$\,pt, MSE $-0.228$), consistent with the intuition that knowing $x_i \bmod m$ for
$m \in \{2,\ldots,101\}$ substantially constrains the set of plausible magnitude values.
Solver accuracy improves by $+7.3$\,pt, primarily because more accurate residue information
refines scoring in Dense and Sieve modes.

\subsection{Scaling Behaviour}
\label{sec:scaling}

Table~\ref{tab:scale_trend} shows how IntSeqBERT's metrics improve with model size.
Magnitude accuracy improves only modestly from Small to Large ($+1.1$\,pt), whereas modulo
accuracy improves substantially ($+10.0$\,pt). This is consistent with the intuition that
modular arithmetic is more compositional and thus benefits more from increased representational
capacity. Solver Top-1 accuracy follows the modulo trend with consistent improvement ($+5.0$\,pt).

\begin{table}[htbp]
\caption{IntSeqBERT scaling: test-split metrics (Small / Middle / Large).}
\label{tab:scale_trend}
\centering
\begin{tabular}{lrrrr}
\hline
Metric              & Small  & Middle & Large  & $\Delta$ (S$\to$L) \\
\hline
Mag Acc (\%)        &  94.73 &  95.71 &  95.85 & $+1.12$ \\
MMA (\%)            &  40.43 &  46.88 &  50.38 & $+9.95$ \\
mod-2 accuracy (\%) &  81.97 &  84.50 &  85.65 & $+3.68$ \\
Solver Top-1 (\%)   &  14.05 &  17.02 &  19.09 & $+5.04$ \\
Magnitude MSE       &  0.228 &  0.164 &  0.142 & $-0.086$\\
\hline
\end{tabular}
\end{table}

The trends in Table~\ref{tab:scale_trend} are visualised in Fig.~\ref{fig:scaling}, which
contrasts the differential scaling of modulo and magnitude accuracy across model sizes.

\begin{figure}[htbp]
  \centering
  \includegraphics[width=\linewidth]{figures/fig3_scaling}
  \caption{Scaling behaviour (Small / Middle / Large). Left: Mag Acc improves only $+1.1$\,pt.
           Centre: MMA for IntSeqBERT improves $+10.1$\,pt (40.3\%$\to$50.4\%).
           Right: Solver Top-1 improves $+5.0$\,pt, with a growing gap over Vanilla.}
  \label{fig:scaling}
\end{figure}

\subsection{Alternative Representations: Digit-Based Tokenisation}
\label{sec:digit-discussion}

A natural alternative to our handcrafted features is digit-based tokenisation, treating each
decimal digit as a separate token, possibly with digit-by-digit generation on the output
side. We have not made an experimental comparison with digit-level representations in this
paper. What our representation targets is not the pattern of decimal digits but the
number-theoretic structure of the integer itself: the modulo stream supplies the residues of
each value under 100 different moduli---information about modular arithmetic---to the model
directly, whereas the same information is only indirectly recoverable from a digit sequence.
Digit-by-digit generation is a reasonable alternative on the output side, but it changes the
problem setting from predicting an integer as a single value to the sequential prediction of
a digit string. A controlled comparison between the two settings is an interesting direction
for future work.

\subsection{Limitations}
\label{sec:limitations}

\textbf{Difficulty with large integers.} Solver accuracy collapses to near zero for
$|x| \geq 10^{20}$ (Huge and Astronomical buckets): CRT-based reconstruction requires
\emph{exact} residues for multiple moduli simultaneously, and errors are frequent enough
at MMA$\approx$50\% to cause failure (Top-1 = 0.09\% in CRT mode).
Approximately 13\% of IntSeqBERT Solver calls return no valid candidate (\texttt{none} mode),
primarily due to CRT failures; approximate CRT or relaxed residue constraints are a natural remedy.

\textbf{Dataset bias.} OEIS sequences are predominantly non-negative (\texttt{nonn}).
The high sign accuracy (98.54\%) may partially reflect this bias.
Similarly, sequences with linear or sub-linear growth remain the majority even after our
filtering, and the high overall magnitude accuracy partially reflects this growth bias;
the contribution of our representation shows up within this bias as the large accuracy gains
in the Large and higher magnitude buckets (Tables~\ref{tab:scale} and~\ref{tab:solver-bucket})
and the effectiveness of the modulo spectrum (Section~\ref{sec:modulo-spectrum}).
Astronomical-bucket metrics are based on few samples and should be interpreted with caution.

\textbf{Compute constraints.} All experiments were conducted on a single consumer GPU
(GeForce RTX 3070 Ti, 8\,GB VRAM), which constrains model size and vocabulary capacity.
Due to this budget, all reported runs use a single random seed (42), and we do not provide
multi-seed variance estimates.
The Vanilla baseline's vocabulary (20,003) is smaller than that of
FACT~\cite{zurich-fact}, which handles values up to several million using larger compute
resources. Consequently, our Vanilla results represent a FACT-equivalent architecture under
our resource constraints rather than a direct reproduction of FACT's reported numbers.
